\documentclass[11pt]{article}
\usepackage{series}

\title{Proto-Spinor Closure and Worldsheet Encoding\\
in Modal Triplet Theory\\
\large A Conditional Local Bridge}
\author{Peter Nero}
\date{July 2026\\Version 5}

\begin{document}
\maketitle

\begin{abstract}
We formulate the proto-spinor/worldsheet relation as a typed local bridge
rather than a derivation of string theory.  Near an aligned q79-compatible
background, a bridge map sends proto-spinor and closure-strain perturbations to
worldsheet coupling perturbations.  If its derivative intertwines the two
quadratic forms and its nonlinear remainder is controlled, the two descriptions
have the same quadratic admissibility test on the selected subspace.  We give
the precise domain, codomain, error estimate, and block conditions.  Circle,
lens, and nil label shared phase, finite transport, and anchoring blocks; they
are not assumed to be literal product factors.  Weyl, Dirac, Majorana, and
twistor languages remain conditional charts.  Worldsheet Weyl invariance,
modular consistency, anomaly cancellation, ghosts, and target-space equations
are independent gates.  The selected q79 trace carrier supplies a concrete
global target for the bridge, but the same-source connection intertwiner is
still required.
\end{abstract}

\section*{Revision note for version 5}
\addcontentsline{toc}{section}{Revision note for version 5}
\begin{description}
\item[Supersedes.] \emph{Proto-Spinor Closure and Worldsheet Encoding in Modal
Triplet Theory: A Conditional Local Bridge}, version~4.
\item[Reason.] Version~4 gives the correct local Hessian-intertwining theorem,
but it did not make the source space, target space, bridge derivative, and
global worldsheet gates sufficiently tangible for a standalone reader.
\item[Resolution.] Version~5 adds a finite-dimensional worked bridge, an
explicit four-object reading map, and a sharper separation between local
quadratic agreement and string-theoretic consistency. It does not replace
the missing normalized q79 bridge derivative with dimensional analogy.
\item[Retained result.] The $C^2$ local bridge theorem, cubic error bound, and
conditional block interpretation remain unchanged.
\item[Remaining boundary.] The same-source q79 derivative, nonflat connection
intertwiner, Weyl and modular consistency, ghosts, anomalies, beta functions,
and a complete physical worldsheet remain open.
\end{description}

\section*{Revision note for version 4}
\addcontentsline{toc}{section}{Revision note for version 4}
\begin{description}
\item[Supersedes.] \emph{Proto-Spinor Closure and Worldsheet Encoding in Modal
Triplet Theory}, version~3.
\item[Reason.] The shadow correspondence lacked one explicit source-to-target
map, a declared validity ball, a nonlinear error estimate, and separation from
full worldsheet consistency.
\item[Resolution.] Version~4 defines a $C^2$ bridge map, proves Hessian
intertwining implies quadratic diagnostic agreement with cubic error, and
lists the independent Weyl, modular, ghost, anomaly, and beta-function gates.
\item[Retained result.] The earlier blockwise matching survives as a
conditional local bridge near alignment.
\item[Remaining boundary.] The normalized q79 bridge derivative and all
global worldsheet and string consistency tests must still be executed.
\end{description}

\tableofcontents

\section{How to Read This Bridge}

This paper does not re-prove the canonical proto-spinor lift.  It assumes a
selected proto-spinor/strain configuration space and asks a different
question: when do its small perturbations have the same local quadratic
admissibility test as perturbations of a chosen worldsheet background?

The bridge has four pieces:
\begin{enumerate}
\item a source configuration $u$ near $u_\ast$;
\item a target worldsheet configuration $c$ near $c_\ast$;
\item a nonlinear map $\mathcal B$ sending source configurations to target
couplings; and
\item a Hessian comparison showing that the derivative
$L=D\mathcal B(u_\ast)$ preserves the leading quadratic diagnostic.
\end{enumerate}
The theorem below is useful only after all four pieces have been constructed.
Matching block dimensions supplies a possible shape for $L$, not its values or
its geometric provenance.

\subsection{A finite-dimensional model}

Let the source and target tangent spaces both be $\mathbb R^2$, take
\[
 H_{\rm ps}=
 \begin{pmatrix}1&0\\0&4\end{pmatrix},
 \qquad
 H_{\rm ws}=I,
 \qquad
 L=
 \begin{pmatrix}1&0\\0&2\end{pmatrix}.
\]
Then
\[
 L^{T}H_{\rm ws}L=H_{\rm ps}.
\]
If $\mathcal B(u)=Lu+r(u)$ with $r(u)=O(\|u\|^2)$ and both diagnostics have
cubic Taylor remainders, their changes agree through quadratic order.  The
difference begins at order $\|u\|^3$.

This toy model contains the whole local mechanism.  The q79 problem is harder
because the spaces are bundles of normalized modes, $L$ must be obtained by
differentiating actual sigma-model couplings, and both Hessians must come from
the same selected source data.

\subsection{Local agreement versus string theory}

A quadratic bridge compares stability and Morse data in a neighborhood of one
background.  It says nothing by itself about large field excursions,
worldsheet renormalization, modular invariance, anomaly cancellation, the
physical state space, or uniqueness of the target description.  Those are
global and quantum consistency questions listed later as independent gates.

\subsection{Argument map}

Sections~2--6 define the two spaces, prove the local bridge, and specialize its
blocks to the selected q79 carrier.  Sections~7--10 separate worldsheet
dimension, spinor/twistor chart language, string consistency, and finite
Standard Model calculations from the local theorem.  Section~11 gives the
actual computation needed to promote the bridge.

\section{Scope}

This paper proves a local equivalence of quadratic diagnostics under explicit
hypotheses.  It does not prove:
\begin{itemize}
\item that a worldsheet is the unique extension of proto-spinor data;
\item that MTT derives string theory, critical dimension, modular invariance,
or a target-space compactification;
\item that a local quadratic bridge extends to the full nonlinear theory; or
\item that a matching number of carrier components identifies two global
bundles.
\end{itemize}

\section{The two configuration spaces}

Let $\mathcal C_{\rm ps}$ be a Sobolev manifold of local proto-spinor,
connection, and strain data near an aligned background $u_\ast$.  A tangent
vector is written
\[
 u=(u_{\rm circ},u_{\rm lens},u_{\rm nil})
 \in E_1\oplus E_2\oplus E_3,
\]
where the labels denote operator/carrier roles.  In the selected q79 interface
their target ranks are $1,2,3$.

Let $\mathcal C_{\rm ws}$ be a Sobolev manifold of worldsheet couplings near a
background $c_\ast$,
\[
 c=(G,B,\Phi,A,\mathcal V,\ldots),
\]
with worldsheet metric coupling $G$, antisymmetric field $B$, dilaton $\Phi$,
gauge connection $A$, and any declared boundary or potential couplings
$\mathcal V$.  The ellipsis is not permission to omit fields needed by the
chosen string model; all active couplings must be listed in an application.

\begin{definition}[Typed bridge map]
A proto-spinor/worldsheet bridge on neighborhoods
$U\subset\mathcal C_{\rm ps}$ and $V\subset\mathcal C_{\rm ws}$ is a $C^2$ map
\[
 \mathcal B:U\longrightarrow V,
 \qquad \mathcal B(u_\ast)=c_\ast.
\]
Its derivative at the background is denoted
$L=D\mathcal B(u_\ast)$.
\end{definition}

This definition replaces an analogy between symbols by a map with a domain,
codomain, and regularity class.

\section{Local remainder control}

Assume $D\mathcal B$ is locally Lipschitz with constant $M_B$.  Taylor's
theorem gives
\[
 \mathcal B(u_\ast+u)
 =c_\ast+Lu+r_B(u),
 \qquad
 \|r_B(u)\|_{\rm ws}\le\frac{M_B}{2}\|u\|_{\rm ps}^2.
\]

Let the two diagnostics have expansions
\begin{align*}
 \mathcal J_{\rm ps}(u_\ast+u)
 &=\mathcal J_{\rm ps}(u_\ast)
 +\frac12\langle u,H_{\rm ps}u\rangle+R_{\rm ps}(u),\\
 \mathcal J_{\rm ws}(c_\ast+v)
 &=\mathcal J_{\rm ws}(c_\ast)
 +\frac12\langle v,H_{\rm ws}v\rangle+R_{\rm ws}(v),
\end{align*}
with cubic bounds on the declared balls.

\begin{theorem}[Conditional quadratic bridge]
Let $E_{\rm sel}\subset T_{u_\ast}\mathcal C_{\rm ps}$ be a closed selected
subspace.  Suppose $L$ is injective on $E_{\rm sel}$ and
\[
 L^*H_{\rm ws}L=H_{\rm ps}\quad\text{on }E_{\rm sel}.
\]
If the displayed Taylor and cubic bounds hold, then there is $C>0$ such that
for sufficiently small $u\in E_{\rm sel}$,
\[
 \left|
 [\mathcal J_{\rm ws}(\mathcal B(u_\ast+u))-\mathcal J_{\rm ws}(c_\ast)]
 -[\mathcal J_{\rm ps}(u_\ast+u)-\mathcal J_{\rm ps}(u_\ast)]
 \right|
 \le C\|u\|_{\rm ps}^3.
\]
\end{theorem}

\begin{proof}
Insert $v=Lu+r_B(u)$ into the worldsheet expansion.  The pure quadratic term
equals the proto-spinor quadratic term by the intertwining identity.  Every
term containing $r_B$ is at least cubic because $r_B=O(\|u\|^2)$, and the two
declared Taylor remainders are cubic.  Their constants combine into $C$.
\end{proof}

\begin{corollary}[Local admissibility equivalence]
If both Hessians are coercive on the selected subspaces, then for sufficiently
small perturbations the bridge preserves strict local admissibility and the
Morse index on $E_{\rm sel}$.
\end{corollary}

This is the rigorous content of the earlier ``shadow bridge.''  It is local,
quadratic to leading order, and conditional on a constructed $\mathcal B$.

The cubic estimate gives a quantitative validity statement rather than exact
nonlinear equivalence.  On a ball of radius $\rho$, the absolute mismatch is
at most $C\rho^3$, while a coercive quadratic signal is of order $\rho^2$.
Thus the relative mismatch is controlled linearly in $\rho$ away from null
directions.  Choosing the ball and proving the constants are part of the
certificate; the notation $O(\|u\|^3)$ is not a license to ignore them.

\section{Blockwise carrier map}

Suppose the selected proto-spinor tangent carrier decomposes as
\[
 E_{\rm sel}=E_1\oplus E_2\oplus E_3,
 \qquad \dim(E_1,E_2,E_3)=(1,2,3),
\]
and choose corresponding worldsheet coupling blocks $W_1,W_2,W_3$.  A
block-preserving derivative has form
\[
 L=L_1\oplus L_2\oplus L_3,
 \qquad L_j:E_j\to W_j.
\]
If
\[
 H_{\rm ps}=\bigoplus_j H_{{\rm ps},j},
 \qquad
 H_{\rm ws}=\bigoplus_j H_{{\rm ws},j},
 \qquad
 L_j^*H_{{\rm ws},j}L_j=H_{{\rm ps},j},
\]
then the bridge theorem holds blockwise.

The roles may be read as follows, provided an application constructs the
actual maps:
\begin{description}
\item[Circle block.] Common $U(1)$ phase/holonomy may map to a compact
worldsheet scalar, Wilson line, or phase coupling.
\item[Lens block.] Finite/projective transport may map to orbifold or twisted
sector data.
\item[Nil block.] Triangular anchoring or boundary termination may map to
boundary couplings, filtration data, or a nilpotent differential.
\end{description}

These are typed possibilities, not a proof that the target worldsheet contains
a literal circle, lens space, or Nil manifold.

\section{The selected q79 target}

The q79 degree-three cover supplies
\[
 \mathcal H_{\rm q79}
 =L_{\rm shared}\otimes
 (\mathcal O\oplus\mathcal A_0\oplus\mathcal A),
 \qquad \operatorname{rank}=1+2+3.
\]
This gives a concrete source for the block ranks.  It does not yet provide the
bridge derivative $L$.  To do so one must construct normalized q79 zero modes
and evaluate how their metric, $B$-field, connection, and scalar variations
enter the worldsheet sigma-model couplings.

The required same-source diagram is
\[
\begin{array}{ccc}
 E_{\rm strain}&\xrightarrow{\mathfrak I_{\rm WW\to q79}}&
 \mathcal H_{\rm q79}\\
 \downarrow &&\downarrow D\mathcal B_{\rm q79}\\
 T\mathcal C_{\rm ps}&\xrightarrow{L}&T\mathcal C_{\rm ws}.
\end{array}
\]
Both routes must agree, and the covariant derivatives and Hessians must
intertwine.  Equality of ranks does not make the square commute.

\section{Why two-dimensionality is not forced}

A two-dimensional worldsheet is a distinguished encoding because local
sigma-model couplings, conformal methods, and string propagation live there.
But no theorem in the proto-spinor carrier alone excludes a worldline, a
higher-dimensional defect, a lattice transfer system, or an operator-algebraic
encoding.  Uniqueness would require a classification theorem with assumptions
strong enough to single out a two-dimensional local conformal theory.

Accordingly, this paper treats the worldsheet as a selected target-compatible
extended chart, not as an inevitable consequence of failed three-dimensional
closure.

\section{Spinorial and twistor chart languages}

The internal proto-spinor is conditional on a lift of rank-three orientation
data.  A spacetime Weyl or Dirac spinor additionally requires a Lorentzian spin
bundle and Dirac operator.  A Majorana chart requires a compatible antilinear
real structure and charge constraints.  A twistor chart requires the relevant
conformal and integrability conditions.

These descriptions may coexist on one coherent slab and be related by typed
maps.  Their coexistence demonstrates representational flexibility, not that
one chart dynamically produces the others.

\section{Worldsheet consistency gates}

Even an exact local quadratic bridge does not define a consistent string
theory.  A selected application must address:
\begin{enumerate}
\item vanishing beta functions or a controlled off-shell worldsheet RG flow;
\item Weyl anomaly and criticality, including ghosts and central charge;
\item modular invariance and the spin-structure/GSO sum where applicable;
\item gauge and gravitational anomaly cancellation;
\item boundary conditions, open/closed consistency, and tadpoles;
\item unitarity and the physical state conditions; and
\item the map from worldsheet beta functions to target-space field equations,
including the order in $\alpha'$ and truncation error.
\end{enumerate}

For heterotic Fu--Yau backgrounds, satisfaction of the Hull--Strominger system
is powerful target-space evidence but does not automatically prove all-order
worldsheet conformal invariance.

\section{Relation to finite SM calculations}

The selected finite-algebra and renormalized-SM packets establish equivalence
at the declared one-shared-physical-primitive/profile standard.  Those packets
do not require a worldsheet derivation and therefore remain valid independently
of this bridge.  Conversely, a worldsheet bridge cannot promote measured
profile input into strict no-knob source selection unless its normalized
overlap integrals and transport scheme emit the values independently.

\section{Execution protocol}

To promote the bridge beyond its current conditional status:
\begin{enumerate}
\item choose one selected q79 Fu--Yau background and its HYM connection;
\item compute normalized harmonic/Dirac representatives for the three carrier
blocks;
\item differentiate the sigma-model coupling map to obtain $L$;
\item compute $H_{\rm ps}$ and $H_{\rm ws}$ from the same action and verify
$L^*H_{\rm ws}L=H_{\rm ps}$ with interval or exact certificates;
\item bound $D^2\mathcal B$ and both cubic remainders on an explicit ball; and
\item run the independent worldsheet consistency gates.
\end{enumerate}

\section{Conclusion}

The proto-spinor/worldsheet relation is now a precise, testable local theorem.
The bridge preserves quadratic admissibility when its derivative intertwines
the Hessians, with a cubic error bound.  The selected q79 carrier gives the
right finite block target, while the actual normalized derivative and global
worldsheet consistency remain to be computed.  This is meaningful progress
without claiming that the existence of a proto-spinor forces string theory.

% BEGIN MTT MANAGED COMPUTATIONAL EVIDENCE
\section*{Computational Evidence and Reproducibility}
The q=79 branch, finite matrix, anomaly, and HYM packets provide concrete compatibility checks for the local proto-spinor/worldsheet encoding. They do not supply the missing global worldsheet, GSO, or physical-bundle theorem. The remaining rows are lower-sector context, and the strict upgrade remains open.

The referenced rows are frozen to the curated results repository state identified below. Hashes are grouped in eight-character blocks for line breaking.
\begin{quote}\small
\textbf{Repository:} \url{https://github.com/PeterNero/mtt-results-repro}\\
\textbf{Commit:} \texttt{31247ebb 5c22f3fb b5443024 365433c6 ee0bff4a}\\
\textbf{Manifest:} \path{release/result_manifest.json}\\
\textbf{Manifest SHA-256:}\\
\texttt{fb399689 60b00584 631dbf53 1a708e18}\\
\texttt{ef928d6b 6d935119 c185d7f6 32b1e7cd}
\end{quote}

Tier labels are quoted verbatim from that manifest. A row used directly supports only the specific computational statement identified above; a corpus-state cross-check does not prove this paper's local theorems; and an open row is evidence of an unresolved obligation, never of closure.

\par\begingroup\dimen0=\pagegoal\advance\dimen0 by -\pagetotal\ifdim\dimen0<7\baselineskip\newpage\fi\endgroup
\paragraph{Rows used directly in this paper.}
\begin{itemize}
\setlength{\itemsep}{0.25em}
\item \path{A22/e6_qpsi_qcd_anomaly} (\emph{derived exact}).\par E6 Qpsi matter/exotic QCD anomaly cancellation audit.
\item \path{A19/hym_wiener_contraction} (\emph{numeric certified}).\par Weighted-theta Fourier-tail and Wiener contraction certificate.
\item \path{A11/q79_exact_audit} (\emph{derived exact}).\par Executable q=79 exact-branch audit.
\item \path{A11/q79_exact_theorem} (\emph{derived exact}).\par CRT q=79 theorem on the selected exact branch.
\item \path{A01/qutrit_weyl_27_matrix} (\emph{derived exact}).\par Sparse 27x27 qutrit-Weyl left-action realization.
\end{itemize}

\par\begingroup\dimen0=\pagegoal\advance\dimen0 by -\pagetotal\ifdim\dimen0<7\baselineskip\newpage\fi\endgroup
\paragraph{Corpus-state cross-checks.}
\begin{itemize}
\setlength{\itemsep}{0.25em}
\item \path{A01/charged_yukawa_higgs_profile} (\emph{profile replay}).\par Versioned Yu, Yd, Ye and lambda\_H profile packet.
\item \path{A14/ckm_prediction_profile} (\emph{numeric certified}).\par Three selected CKM profile rows and uncertainty comparison.
\item \path{A01/current_global_lock} (\emph{profile replay}).\par Current non-looping global status and source-certificate map.
\item \path{A01/direct_k_higgs_row} (\emph{derived exact}).\par Promoted direct K\_threshold.Omega\_H.lambda row.
\item \path{A04/final_12_of_12_audit} (\emph{profile replay}).\par Twelve-obligation embedded renormalized-SM equivalence audit.
\item \path{A13/gr_tt_support} (\emph{derived exact}).\par Exact-branch internal TT support certificate; physical normalization remains open.
\item \path{A40/neutral_two_primitive_profile} (\emph{profile replay}).\par Measured two-splitting neutrino profile, masses and Dirac Yukawa rows.
\item \path{A02/precision_15_source_transport} (\emph{profile replay}).\par Fifteen measured source coordinates, Jacobian and covariance transport.
\item \path{A02/precision_8x8_workspace} (\emph{profile replay}).\par Eight-coordinate SMDR output with positive-definite 8x8 covariance.
\item \path{A01/strict_pew_row} (\emph{derived exact}).\par Promoted P\_EW source row at the declared one-shared-primitive standard.
\end{itemize}

\par\begingroup\dimen0=\pagegoal\advance\dimen0 by -\pagetotal\ifdim\dimen0<7\baselineskip\newpage\fi\endgroup
\paragraph{Open boundary (not evidence of closure).}
\begin{itemize}
\setlength{\itemsep}{0.25em}
\item \path{A05/strict_upgrade_ledger} (\emph{open}).\par Current 2/9 strict no-knob upgrade ledger.
\end{itemize}

No imported row changes theorem ownership or promotes a neighboring claim: all local statements retain their stated hypotheses, domains, and limitations.
% END MTT MANAGED COMPUTATIONAL EVIDENCE

\begin{thebibliography}{9}
\bibitem{Polchinski}
J.~Polchinski, \emph{String Theory}, Vols. 1--2, Cambridge University Press,
1998.

\bibitem{FuYau}
J.-X.~Fu and S.-T.~Yau,
\emph{The theory of superstring with flux on non-Kahler manifolds and the
complex Monge--Ampere equation}, J. Differential Geom. 78 (2008).

\bibitem{MTTBridge}
P.~Nero, \emph{Selected q79 Trace-Split CLN Carrier and World-in-World
Bridge}, internal theorem packet, 2026.
\end{thebibliography}



\end{document}
